\documentclass{article}
\usepackage{amsmath,amssymb,amsthm,thmtools,enumitem,bm,graphicx,xcolor,microtype}
\usepackage[T1]{fontenc}

\usepackage[style = ext-numeric, sorting = nyt, giveninits = false, maxbibnames = 8, minbibnames = 3, date = year, isbn = false, doi = true, eprint = true, url = true, backref = true]{biblatex}
\addbibresource{../../ZoteroDatabase.bib}
\renewbibmacro{in:}{}
\DeclareFieldFormat{titlecase}{\MakeSentenceCase*{#1}}
\DeclareFieldFormat{titlecase:journaltitle}{{#1}}

\usepackage[unicode]{hyperref}

\theoremstyle{definition}
\newtheorem{dfn}{Definition}[section]
\newtheorem{ntt}[dfn]{Notation}
\newtheorem{ex}[dfn]{Example}
\newtheorem{case}{Case}
\newtheorem{cnd}{Condition}

\theoremstyle{remark}
\newtheorem{rmk}[dfn]{Remark}
\newtheorem*{sroof}{Sketch of Proof}

\theoremstyle{plain}
\newtheorem{thm}[dfn]{Theorem}
\newtheorem*{thm*}{Theorem}
\newtheorem{prop}[dfn]{Proposition}
\newtheorem*{prop*}{Proposition}
\newtheorem{lem}[dfn]{Lemma}
\newtheorem*{lem*}{Lemma}
\newtheorem{cor}[dfn]{Corollary}
\newtheorem{fct}[dfn]{Fact}
\newtheorem{conj}[dfn]{Conjecture}
\newtheorem{qst}[dfn]{Question}
\newtheorem{clm}{Claim}

\renewcommand{\qedsymbol}{$\blacksquare$}

\declaretheoremstyle[headfont = \normalfont\itshape, qed = $\blacksquare_{\mathrm{Claim}}$]{prooc}
\declaretheorem[style = prooc, name = Proof of Claim, numbered = no]{prooc}

\def\sectionautorefname~#1\null{Chapter~#1\null}
\def\subsectionautorefname~#1\null{Section~#1\null}

\title{}
\author{}
\date{}

\begin{document}

% \maketitle

% \tableofcontents

\paragraph{Introduction}

We will compute the dp-rank of ordered abelian groups (OAGs) in the pure language $\mathcal{L} = \left\{ <,+ \right\}$.
This is NOT a new result.
Gurevich and Schmitt (\cite{Gurevich1984TheOrd}), showed that all OAGs have NIP, and NIP is equivalent to having dp-rank $\leq |\mathcal{L}| + \aleph_0$.
Furthermore, Farré \cite{Farre2017StrOrd}, Dolich and Goodrick \cite{Dolich2018ChaStr}, and Halevi and Hasson \cite{Halevi2019StrDep} independently established that an OAG has finite dp-rank if and only if it has finite spines and there are only finitely many infinite ribs.
Farré also gave a formula for the dp-rank of an OAG.
The purpose of this note is to practice AI-assisted proof writing.

\paragraph{Dp-rank}

Let $T$ be a (complete) theory in a language $\mathcal{L}$.
We take elements from its monster model.
The \emph{dp-rank} of $T$ is at least a cardinal $\kappa$ if there are formulas $(\varphi_{\alpha}(x,\overline{y}_{\alpha}))_{\alpha < \kappa}$, parameters $(\overline{b}_{\alpha, k})_{\alpha < \kappa, k < \omega}$, and elements $(a_{\eta})_{\eta}$ where $\eta
$ ranges over the functions $\kappa \to \omega$, such that
\begin{equation*}
	\models \varphi_{\alpha}(a_{\eta}, \overline{b}_{\alpha, k}) \iff \eta(\alpha) = k
\end{equation*}
for all $\alpha < \kappa$, $k < \omega$, and $\eta: \kappa \to \omega$ (such a configuration is called an ICT-pattern of depth $\kappa$).
Note that it is possible for a theory to have dp-rank greater than all positive integers and less than $\aleph_0$.
NIP is known to be equivalent to having dp-rank at most $|\mathcal{L}| + \aleph_0$, as we mentioned earlier.

There is a useful criterion for evaluating dp-rank, which may be proved in the same way as \cite[Fact 2.10]{Dolich2011DpmBas}.
Let $n$ be a positive integer.
Suppose that for any $I \models \mathrm{DLO}$, any indiscernible sequence $(\overline{a}_i)_{i \in I}$, and any element $b$, there is a decomposition $I_0 < J_0 < I_1 < \dots < J_{n - 1} < I_n$ of $I$ such that
\begin{itemize}
	\item $I_k$ is without endpoints for all $0 \leq k \leq n$.
	\item $|J_k| \leq 1$ for all $0 \leq k < n$.
	\item $\mathrm{tp}(\overline{a}_i / b) = \mathrm{tp}(\overline{a}_j / b)$ for all $0 \leq k \leq n$ and $i, j \in I_k$.
\end{itemize}
Then the dp-rank of $T$ is at most $n$.

\paragraph{Model theory of OAGs}

OAGs are well-studied in model theory.
Although Schmitt \cite{Schmitt1982ModThe} gave a quite detailed analysis early in 1980s, we use Cluckers and Halpuczok \cite{Cluckers2011QuaEli} as a primary reference because it is modern and easy to access.
A paper by Koki Okura (the author of this note) (\cite{Okura2026DisOrd}) is also a useful reference, we hope.
We leave a detailed account to the references and only give a brief overview here.

Fix an OAG $G$.
For $n \geq 2$ and $a \in G$, convex subgroups $\mathfrak{s}_n(a)$ and $\mathfrak{t}_n(a)$ are defined uniformly in $a$.
They behave as valuations on $G$ (technically they do not satisfy the valuation axioms, but they are close enough).
Hence, we can define \emph{spines} $\mathcal{S}_n$ and $\mathcal{T}_n$, and \emph{ribs} $R_H^{\mathfrak{s}_n}$ and $R_H^{\mathfrak{t}_n}$, which are counterparts of value groups and residue fields of valued fields.
These are essential for the relative quantifier elimination given by Cluckers and Halpuczok, and hence for the model theory of OAGs.

\paragraph{Expected result}

Since every OAG has NIP (\cite{Gurevich1984TheOrd}), it has dp-rank at most $\aleph_0$.
We will show the followings:
Fix an OAG $G$.
First, if $G$ has an infinite spine (that is, $\mathcal{S}_p$ is infinite for some prime $p$), then it has dp-rank $\aleph_0$.
Suppose that $G$ has finite spines.
For each prime $p$, let $\mathcal{S}_p^{\infty} = \left\{ H \in \mathcal{S}_p \setminus \left\{ \emptyset \right\} \mid \left| R_H^{\mathfrak{s}_p} \right| = \infty \right\}$.
Then the dp-rank of $G$ equals $1 + \sum_{p: \text{prime}} \left| \mathcal{S}_p^{\infty} \right|$ (if the right-hand side is infinite, then the dp-rank is $\aleph_0$).

\paragraph{First theorem: infinite spine case}

We fix an OAG $G$ hereafter.
We will prove the following theorem:

\begin{thm*}
	If there is a prime $p$ such that $\mathcal{S}_p$ is infinite, then the dp-rank of $G$ is $\aleph_0$.
\end{thm*}

We sketch the proof.
As we know that the dp-rank is at most $\aleph_0$, it suffices to show that the dp-rank is at least $\aleph_0$.
We may assume that $G$ is a monster model.
Fix a prime $p$ such that $\mathcal{S}_p$ is infinite.
We can choose an infinite sequence $H_0 < H_1 < \dots$ from $\mathcal{S}_p \setminus \left\{ \emptyset \right\}$.
For each $k < \omega$, there is an element $a_k \in G$ such that $\mathfrak{s}_p(a_k) = H_k$ and $\mathfrak{s}_{p^r}(a_k) < H_{k+1}$ for all $r \geq 1$, since there is $H'_k \in \mathcal{T}_p$ such that $H_k \leq H'_k < H_{k+1}$, and $\mathfrak{s}_p(a) \leq \mathfrak{s}_{p^2}(a) \leq \dots \leq \mathfrak{t}_p(a)$ for all $a \in G$.
Then for any $1 \leq r \leq N + 1 < \omega$, we have $\mathfrak{s}_{p^r}(a_0 + p a_1 + \dots + p^N a_N) = H_{r-1}$.
Now we choose $H_{0,0} < H_{0,1} < \dots < H_{1,0} < H_{1,1} < \dots$ from $\mathcal{S}_p$.
For any function $\eta: \omega \to \omega$, we can choose $a_{\eta} \in G$ such that $\mathfrak{s}_{p^r}(a_{\eta}) = H_{r-1, \eta(r-1)}$ for all $r \geq 1$.
These give an ICT-pattern of depth $\aleph_0$, and hence the dp-rank of $G$ is at least $\aleph_0$.

\paragraph{Second theorem: lower bound}
We assume that $G$ is nontrivial and has finite spines hereafter.
Let $\mathbb{P}$ be the set of primes.
We will prove the following theorem:

\begin{thm*}
	The dp-rank of $G$ is at least $1 + \sum_{p \in \mathbb{P}} \left| \mathcal{S}_p^{\infty} \right|$ (if the right-hand side is infinite, then the dp-rank is $\aleph_0$).
\end{thm*}

The sketch of the proof is as follows.
We may assume that $G$ is a monster model.
Let $P_0$ be a finite set of primes, and for each $p \in P_0$, let $H_{p,0} < H_{p,1} < \dots < H_{p,n_p-1}$ be finitely many elements of $\mathcal{S}_p^{\infty}$.
For each $p \in P_0$ and $0 \leq i < n_p$, we can choose $(a_{p,i}^k)_{k < \omega}$ such that $\mathfrak{s}_p(a_{p,i}^k) = H_{p,i}$ and (if $i < n_p - 1$) $\mathfrak{s}_{p^r}(a_{p,i}^k) < H_{p,i+1}$ for all $k < \omega$ and $r \geq 1$, and the images of $(a_{p,i}^k)_{k < \omega}$ in the rib $R_{H_{p,i}}^{\mathfrak{s}_p}$ are pairwise distinct.
For each $\zeta: n_p \to \omega$, let $b_{p,\zeta} = a_{p,0}^{\zeta(0)} + p a_{p,1}^{\zeta(1)} + \dots + p^{n_p - 1} a_{p,n_p - 1}^{\zeta(n_p - 1)}$.
Then, for each $0 \leq i < n_p$, we have $\mathfrak{s}_{p^{i+1}}(b_{p,\zeta}) = H_{p,i}$, and the image of $b_{p,\zeta}$ in $R_{H_{p,i}}^{\mathfrak{s}_{p^{i+1}}}$ is the same as that of $p^i a_{p,i}^{\zeta(i)}$.
We can construct an ICT-pattern of depth $n_p$ with these elements.
Furthermore, by the Chinese remainder theorem (\cite[Lemma 2.1 (2)]{Okura2026DisOrd}), we can construct an ICT-pattern of depth $\sum_{p \in P_0} n_p$.
Finally, let $N = \prod_{p \in P_0} p^{n_p}$.
Since $N G$ is unbounded in $G$ and translation by an element of $N G$ does not affect any of the congruence rows constructed above, we can construct intervals $(c_0, d_0) < (c_1, d_1) < \dots$ such that every congruence path has a realization in each interval.
This leads to an ICT-pattern of depth $1 + \sum_{p \in P_0} n_p$.
A Compactness argument leads to the desired result.

\paragraph{Third theorem: upper bound}

We further assume that $\sum_{p \in \mathbb{P}} \left| \mathcal{S}_p^{\infty} \right|$ is finite.

\begin{thm*}
	The dp-rank of $G$ is at most $1 + \sum_{p \in \mathbb{P}} \left| \mathcal{S}_p^{\infty} \right|$.
\end{thm*}

We make use of the relative quantifier elimination given by Cluckers and Halpuczok (\cite{Cluckers2011QuaEli}) and its consequence.
\cite[Proposition 2.26]{Okura2026DisOrd} is specialized as follows:

\begin{prop*}
	Let $G$ be an OAG with finite spines, and let $f: G \to G$ be a partial function.
	Suppose that $f$ satisfies the following conditions:
	\begin{itemize}
		\item $\mathrm{Dom}(f)$ is a subgroup of $G$.
		\item For every $a \in \mathrm{Dom}(f)$, the values of $a$ and $f(a)$ with respect to the valuations $\mathfrak{s}_{p^r}$ and $\mathfrak{t}_p$ coincide for all $p \in \mathbb{P}$ and $r \geq 1$.
		\item $f$ is an ordered group embedding that preserves the predicates $x =_{\bullet} k$ and $x \equiv_{p^r, \bullet} k$ for all $k$, $p$, and $r$.
	\end{itemize}
	Then $f$ is elementary.
\end{prop*}

(Note that the predicates $D_{p^r}^{[p^s]}$ are trivial in finite spines case, because for each $H \in \mathcal{S}_p \setminus \left\{ \emptyset \right\}$, if it has a successor $H'$ in $\mathcal{S}_p$then $H^{[p^s]} = H' + p^s G$, and if $H$ has no successor then $H^{[p^s]} = G$.)

Let $I \models \mathrm{DLO}$, $(\overline{a}_i)_{i \in I}$ be an indiscernible sequence, and $b \in G$.
For each $i \in I$, let $M_i$ be the subgroup generated by $\overline{a}_i$.
$(M_i)_{i \in I}$ is also an indiscernible sequence.
We first prove the following lemma:

\begin{lem*}
	Suppose that there is a decomposition $I_0 < J < I_1$ of $I$ such that $I_0$ and $I_1$ are without endpoints, $|J| \leq 1$, and there are a sequence $(c_i)_{i \in I}$ from $(M_i)_{i \in I}$ and an nonzero integer $m$ such that one of the following holds:
	\begin{itemize}
		\item The sign of $(c_i - mb)_{i \in I_0}$ and that of $(c_i - mb)_{i \in I_1}$ are opposite, or
		\item for some $k \in \mathbb{Z}$, $p \in \mathbb{P}$, and $H \in \mathcal{T}_p \setminus \left\{ \emptyset \right\}$, the sign of $((c_i - mb) / H - k_H)_{i \in I_0}$ and that of $((c_i - mb) / H - k_H)_{i \in I_1}$ are opposite.
	\end{itemize}
	Suppose further that there is another such decomposition $I'_0 < J' < I'_1$, $(c'_i)_{i \in I}$, and $m'$ such that $I'_0 \cap I_1 \neq \emptyset$.
	Then a contradiction arises.
\end{lem*}

For instance, suppose that $(c_i - b)_{i \in I_0}$ and $(c_i - b)_{i \in I'_0}$ are negative and that $(c'_i - b)_{i \in I_1}$ and $(c'_i - b)_{i \in I'_1}$ are positive.
If we choose $i < j$ from $I'_0 \cap I_1$, then we have $c'_j < c_i$.
However, if $i \in I_0$, $j \in I'_1$, and $i < j$, then we have $c_i < c'_j$.
These contradict the indiscernibility of $(M_i)_{i \in I}$.

By the lemma, we may choose a decomposition $I_0 < J < I_1$ of $I$ such that for any $i, j \in I_0$ (or $I_1$), a sequence $(c_i)_{i \in I}$ from $(M_i)_{i \in I}$, and an nonzero integer $m$, the following holds:
\begin{itemize}
	\item The sign of $c_i - mb$ and that of $c_j - mb$ are the same.
	\item For every $p \in \mathbb{P}$, $\mathfrak{t}_p(c_i - b) = \mathfrak{t}_p(c_j - b)$.
	\item For every $p \in \mathbb{P}$ and $H \in \mathcal{T}_p \setminus \left\{ \emptyset \right\}$, if $G / H$ is discrete, for every $k \in \mathbb{Z} \setminus \left\{ 0 \right\}$, the sign of $(c_i - mb) / H - k_H$ and that of $(c_j - mb) / H - k_H$ are the same.
\end{itemize}
If we define $f: \langle M_i, b \rangle \to \langle M_j, b \rangle$ ($\langle M_i, b \rangle$ is the generated subgroup) by $f(c_i - mb) = c_j - mb$, this ensures that $f$ satisfies the conditions of the proposition about $\mathfrak{t}_p$, ordered group embeddingness, and $x =_{\bullet} k$.

We also need the following lemma:

\begin{lem*}
	Let $p \in \mathbb{P}$ and $H \in \mathcal{S}_p^{\infty}$.
	Suppose that there is $i_0 \in I$, and there are a sequence $(c_i)_{i \in I}$ from $(M_i)_{i \in I}$, $r \geq 1$, and an nonzero integer $m$ such that $\mathfrak{s}_{p^r}(c_{i_0} - mb) < H$ and $\mathfrak{s}_{p^r}(c_i - mb) = H$ for all $i \in I \setminus \left\{ i_0 \right\}$.
	Further suppose that there are such $i'_0 \in I$, $(c'_i)_{i \in I}$, $r'$, and $m'$ such that $i_0 \neq i'_0$.
	Then a contradiction arises.
\end{lem*}

We assume that $i_0 < i'_0$.
Also we may assume that $r = r'$ by multiplying a suitable power of $p$, and that $m$ and $m'$ are powers of $p$ by multiplying a suitable integers coprime to $p$.
For instance, let $m = 1$ and $m' = p$.
$c_{i_0} \equiv b$ and $c'_{i'_0} \equiv pb$ modulo $H + p^r G$, and then $s(c'_{i'_0} - c_{i_0}) = b$ modulo $H + p^r G$ for $s$ which is an inverse of $p - 1$ in $\mathbb{Z} / p^r \mathbb{Z}$.
Then $s c'_{i'_0} - (s-1) c_{i_0} \in H + p^r G$, or equivalently $\mathfrak{s}_{p^r}(s c'_{i'_0} - (s-1) c_{i_0}) < H$.
Then we can show that $\mathfrak{s}_{p^r}(c'_j - c'_{i'_0}) < H$ for all $j \in I_{>i'_0}$, by indiscernibility.
Hence $\mathfrak{s}_{p^r}(c'_j - pb) < H$ for all $j \in I_{>i'_0}$, a contradiction.

For $p \in \mathbb{P}$ and $r \geq 1$, let $H = \max \left\{ \mathsf{s}_{p^r}(c_i - mb) \mid i \in I \right\}$.
If $(\mathsf{s}_{p^r}(c_i - mb))_{i \in I}$ is not constant, then there is only one $i_0 \in I$ such that $\mathsf{s}_{p^r}(c_{i_0} - mb) < H$.
Combining this observation with the lemma, the following holds:
We can choose $\sum_{p \in \mathbb{P}} \left| \mathcal{S}_p^{\infty} \right|$ or less points in $I$ such that, for any convex subset $C$ of $I$ containing no point chosen it holds:
For any $i,j \in C$, any sequence $(c_i)_{i \in I}$ from $(M_i)_{i \in I}$, $p \in \mathbb{P}$, $r \geq 1$, and any nonzero integer $k, m$, we have $\mathfrak{s}_{p^r}(c_i - mb) = \mathfrak{s}_{p^r}(c_j - mb)$ and $c_i - mb \equiv_{p^r, \bullet} k \iff c_j - mb \equiv_{p^r, \bullet} k$ (we did not mention the predicates $x \equiv_{p^r, \bullet} k$ in the lemma, but a similar reasoning should work).
It ensures that $f$ defined above preserves $\mathfrak{s}_{p^r}$ and $x \equiv_{p^r, \bullet} k$.

Now we may apply the criterion mentioned in the first page to show the theorem.

\paragraph{Instruction of Introduction}

Add a concise introduction. Include the following contents:
\begin{itemize}
	\item The theorems proven are far from original. Those were proven by Farré, and relevant results were given by Dolich and Goodrick, and Halevi and Hasson.
	\item The purpose of this note is to practice AI-assisted proof writing.
	\item The author gave sketches of the proofs, AI (Chat GPT 5.6 Sol and Grok 4.6) completed the proofs, and the author reviewd and revised the proofs.
	\item Notation follows the paper by the author.
\end{itemize}


%\addcontentsline{toc}{section}{References}
\begin{sloppypar}
	\printbibliography
\end{sloppypar}

\end{document}